\documentclass[12pt]{article}
\usepackage[utf8]{inputenc}
\usepackage[T1]{fontenc}
\usepackage{hyperref}
\title{Sab\"um Language Resources for K-12 Curriculum Development}
\author{Compiled by Claude Sonnet 5.5 (claude-sonnet-5-5) for Andr\'as Kornai}
\date{}
\begin{document}
\maketitle

\section{Introduction}
Sab\"um (ISO 639-3: sbo) is an Aslian language of the Austroasiatic family, belonging to the Central Aslian (Senoic) group. It was spoken in the Lenggong area of Perak, Malaysia. It is dormant: no first-language speakers are known, and sources differ on when it fell out of use (Wikipedia says nearly extinct by 1976, ASJP gives 1980). Its speakers are reported to have shifted to Temiar. It is closely related to Lanoh and Semnam, and the Malaysian official label ``Lanoh'' covers several Central Aslian varieties including Sab\"um. Searches found no Sab\"um-specific textbook, dictionary, grammar or school programme. Only short reference entries and a small word sample exist.

\section{Primary Educational Resources}
No primers, textbooks or school materials in Sab\"um were found. The language is not known to be taught in schools.

\section{Linguistic and Reference Materials}
\begin{itemize}
\item Wikipedia, Sab\"um language (classification, location, a short comparative word sample with Semnam): \url{https://en.wikipedia.org/wiki/Sab%C3%BCm_language}
\item ASJP database entry for Sabum (wordlist-based record): \url{https://asjp.clld.org/languages/SABUM}
\item Endangered Languages Project entry (status: dormant, no known L1 speakers): \url{https://endangeredlanguages.com/lang/5358?hl=en}
\item Wikipedia, Senoic languages (context for the Central Aslian group): \url{https://en.wikipedia.org/wiki/Senoic_languages}
\item Wikipedia, Lanoh language (closest documented relative): \url{https://en.wikipedia.org/wiki/Lanoh_language}
\end{itemize}

\section{Religious and Community Resources}
No Sab\"um religious materials were found. The Lanoh people article gives community background for the related group: \url{https://en.wikipedia.org/wiki/Lanoh_people}

\section{Digital and Audio Resources}
No Sab\"um audio was found. The DoBeS ``Tongues of the Semang'' project documented the related Lanoh and Semnam varieties, and its archive may hold relevant comparative material: \url{https://dobes.mpi.nl/?p=1273}

\section{Curriculum Development Context}
Because almost nothing is documented for Sab\"um itself, a revitalization curriculum would have to be built from a small word sample plus comparative work with Lanoh/Semnam, and any reconstruction should be done with the descendant community and linguists. This dossier does not claim such materials exist. The community's current language (Temiar, Malay) would be the practical medium of instruction.
\end{document}